% Lösungen Einheit 3 von 3 – Quadratwurzeln (zusammenbau v0.1)
\einheitenkopf{Einheit 3 von 3 · Quadratwurzeln}

\erg{21}{a) erst das Quadrat – es steht unter der Wurzel}
%% TODO Lösung der Erklärzeile 22g) fehlt in der Bank
\erg{22}{a) ja, Quadratzahl – elf mal elf \quad b) $7$, Probe: $7 \cdot 7 = 49$ \quad c) $8$, Probe: $8 \cdot 8 = 64$ \quad d) $10$, Probe: $10 \cdot 10 = 100$ \quad e) $15$, Probe: $15 \cdot 15 = 225$ \quad f) $20$, Probe: $20 \cdot 20 = 400$ \quad g) \ldots \quad h) $\sqrt{11} \approx 3{,}32$ \quad i) zwischen $7$ und $8$, denn $49 < 60 < 64$; Taschenrechner: $\sqrt{60} \approx 7{,}75$ \quad j) $\sqrt{3} \approx 1{,}732 > 1{,}7$ \quad k) $-2{,}7 < \frac{5}{2} < 2{,}6 < \sqrt{7}$, denn $\sqrt{7} \approx 2{,}65$ \quad l) $13$ \quad m) $(-8)^2 = 64$, $\sqrt{64} = 8$ \quad n) $-\sqrt{49} = -7$; $\sqrt{-49}$ hat keinen Wert, weil kein Quadrat negativ ist \quad o) $0{,}7$, denn $0{,}7 \cdot 0{,}7 = 0{,}49$ \quad p) $\sqrt{6{,}25} = 2{,}5$, die Seite ist $2{,}5$ m lang \quad q) $\sqrt{49 + 576} = \sqrt{625} = 25$ \quad r) $4$, denn $4 \cdot 4 \cdot 4 = 64$ \quad s) $\sqrt{(-7)^2} = 7$, denn erst das Quadrat: $(-7)^2 = 49$, und die Wurzel ist nie negativ}
\erg{23}{a) Die Wurzel halbiert statt gezogen; richtig: $\sqrt{64} = 8$, denn $8 \cdot 8 = 64$ \quad b) Die Wurzel ist die nicht negative Zahl, die mit sich selbst malgenommen $49$ ergibt – also $\sqrt{49} = 7$.}
